% ============================================================================
% SPANISCH LANGENTWURF - DS 9: Vocabulario y Gramática (Einstimmung)
% ============================================================================
% Kurs: Jahrgangsübergreifend Spanisch Spätbeginner (Kl. 10, 11, 12)
% Datum: Mo 02.02.2026
% Besonderheit: TAG VOR DER BEOBACHTUNG - Wiederholung ohne neue Inhalte
% ============================================================================

\documentclass[11pt,a4paper]{article}

% ============================================================================
% PACKAGES
% ============================================================================
\usepackage[utf8]{inputenc}
\usepackage[T1]{fontenc}
\usepackage[ngerman]{babel}
\usepackage{geometry}
\usepackage{fancyhdr}
\usepackage{lastpage}
\usepackage[table]{xcolor}
\usepackage{tcolorbox}
\usepackage{enumitem}
\usepackage{tabularx}
\usepackage{longtable}
\usepackage{array}
\usepackage{booktabs}
\usepackage{hyperref}
\usepackage{ifthen}
\usepackage{amssymb}

% ============================================================================
% KONFIGURATION
% ============================================================================

% --- TIER ---
\newcommand{\plantier}{1}

% --- FACH ---
\newcommand{\fachid}{spanisch}
\newcommand{\fach}{Spanisch}
\newcommand{\fachfarbe}{c41e3a}

% ============================================================================
% VARIABLEN
% ============================================================================

% --- Metadaten ---
\newcommand{\jahrgangsstufe}{10/11/12 (jahrgangsübergreifend)}
\newcommand{\schulart}{Gesamtschule}
\newcommand{\stundenthema}{Vocabulario y Gramática -- Einstimmungsstunde}
\newcommand{\reihenthema}{Beobachtungsstunde Februar 2026}
\newcommand{\stundennummer}{9}
\newcommand{\reihenstunden}{10}
\newcommand{\unterrichtsdatum}{02.02.2026}
\newcommand{\unterrichtszeit}{[Lt. Stundenplan]}
\newcommand{\raum}{[Fachraum]}

% --- Lehrkraft ---
\newcommand{\lehrkraft}{[Name]}
\newcommand{\ausbildungsstand}{Lehrkraft}

% --- Lerngruppe ---
\newcommand{\lerngruppengroesse}{10}
\newcommand{\maennlich}{1}
\newcommand{\weiblich}{9}
\newcommand{\divers}{0}

% --- Kompetenzen/Niveau ---
\newcommand{\gerniveau}{A2--B1}
\newcommand{\lehrwerk}{A\_tope.com}
\newcommand{\lehrwerkseite}{--}

% ============================================================================
% LAYOUT
% ============================================================================
\geometry{left=2.5cm, right=2.5cm, top=2.5cm, bottom=2.5cm}
\definecolor{fach}{HTML}{\fachfarbe}
\definecolor{gray}{HTML}{666666}
\definecolor{lightgray}{HTML}{f5f5f5}
\definecolor{successgreen}{HTML}{2a7a4b}
\definecolor{warningorange}{HTML}{b35c00}
\definecolor{gruppeA}{HTML}{2c5aa0}
\definecolor{gruppeB}{HTML}{c41e3a}
\definecolor{overlap}{HTML}{6a0dad}

\tcbuselibrary{skins,breakable}

% Farbige Boxen
\newtcolorbox{infobox}[1][]{
    colback=lightgray,
    colframe=fach,
    fonttitle=\bfseries,
    title=#1,
    breakable
}

\newtcolorbox{scorebox}{
    colback=white,
    colframe=fach,
    fonttitle=\bfseries,
    title=Didaktik-Score,
    breakable
}

\newtcolorbox{onlinebox}{
    colback=white,
    colframe=warningorange,
    fonttitle=\bfseries,
    title={ONLINE-VARIANTE -- Bei Lehrerausfall},
    breakable
}

\newtcolorbox{gamebox}{
    colback=white,
    colframe=overlap,
    fonttitle=\bfseries,
    title={GAMIFICATION-ELEMENT},
    breakable
}

% Kopf-/Fußzeile
\pagestyle{fancy}
\fancyhf{}
\fancyhead[L]{\small\textcolor{gray}{\fach{} -- Jg. \jahrgangsstufe{} -- \stundenthema}}
\fancyhead[R]{\small\textcolor{gray}{Langentwurf Tier 1}}
\fancyfoot[C]{\small\thepage{} / \pageref{LastPage}}
\renewcommand{\headrulewidth}{0.4pt}
\renewcommand{\footrulewidth}{0pt}

% ============================================================================
% DOKUMENT
% ============================================================================
\begin{document}

% ============================================================================
% DECKBLATT
% ============================================================================
\begin{titlepage}
    \centering
    \vspace*{2cm}

    {\large\textcolor{gray}{\schulart}}\\[2cm]

    {\Huge\textcolor{fach}{\textbf{Langentwurf}}}\\[0.5cm]
    {\large Tier 1 -- Vollständig}\\[2cm]

    {\LARGE\textbf{\stundenthema}}\\[0.5cm]
    {\large im Rahmen der Unterrichtsreihe}\\[0.3cm]
    {\Large\textit{\reihenthema}}\\[1cm]

    \textcolor{overlap}{\Large\textbf{--- TAG VOR DER BEOBACHTUNG ---}}\\[0.3cm]
    {\normalsize Aktivierung und Selbstvertrauen stärken}\\[0.3cm]
    {\normalsize \textbf{KEINE neuen Inhalte -- Bekanntes in neuer Kombination}}\\[1.5cm]

    \begin{tabular}{rl}
        \textbf{Fach:} & \fach{} (\gerniveau) \\[0.3cm]
        \textbf{Klasse:} & \jahrgangsstufe \\[0.3cm]
        \textbf{Lehrwerk:} & \lehrwerk \\[0.3cm]
        \textbf{Datum:} & \underline{\hspace{1cm}\unterrichtsdatum\hspace{1cm}} (Montag) \\[0.3cm]
        \textbf{Zeit:} & \underline{\hspace{3cm}} (Doppelstunde) \\[0.3cm]
        \textbf{Raum:} & \underline{\hspace{3cm}} \\[0.3cm]
    \end{tabular}

    \vfill

    {\large Lehrkraft: \lehrkraft}\\[0.3cm]
    {\normalsize\textcolor{gray}{\ausbildungsstand}}\\[1cm]

\end{titlepage}

% ============================================================================
% INHALTSVERZEICHNIS
% ============================================================================
\tableofcontents
\newpage

% ============================================================================
% 1. BEDINGUNGSANALYSE
% ============================================================================
\section{Bedingungsanalyse}

\subsection{Lerngruppe}
\begin{tabular}{ll}
    Kursgröße: & \lerngruppengroesse{} SuS (\maennlich{} m / \weiblich{} w / \divers{} d) \\
    Jahrgangsstufe: & Kl. 10, 11, 12 (jahrgangsübergreifend) \\
    Sprachniveau: & \gerniveau{} (GER) -- heterogen \\
    Fremdsprache: & 3. Fremdsprache (Spätbeginner) \\
\end{tabular}

\subsection{Schülerprofile und Gruppenzuordnung}

\textbf{Besonderheit dieser Stunde:} Tag vor der Beobachtung -- Fokus auf Routine und Sicherheit!

\vspace{0.5em}

\begin{tabular}{|l|c|l|c|p{4.5cm}|}
\hline
\textbf{Name} & \textbf{Kl.} & \textbf{Gruppe} & \textbf{NP} & \textbf{Profil} \\
\hline
\textcolor{gruppeA}{E.H.} & 10 & \textcolor{gruppeA}{A (A2)} & 15 & TOP, Peer-Tutor, Franz. Hintergrund \\
\hline
\textcolor{gruppeA}{L.N.} & 10 & \textcolor{gruppeA}{A (A2)} & 13 & Proaktiv, manchmal etwas faul \\
\hline
\textcolor{gruppeA}{C.P.} & 10 & \textcolor{gruppeA}{A (A2)} & 12 & Ruhig, gute Aussprache \\
\hline
\textcolor{gruppeA}{V.S.} & 10 & \textcolor{gruppeA}{A (A2)} & 11 & Konstant, fleißig, ruhig \\
\hline
\textcolor{gruppeA}{L.K.} & 11 & \textcolor{gruppeA}{A (A2)} & 10 & Kl.11 Nachholer, manchmal abgelenkt \\
\hline
\textcolor{gruppeB}{H.P.} & 11 & \textcolor{gruppeB}{B (B1)} & 10 & Späteinsteiger, Grundlagenlücken \\
\hline
\textcolor{gruppeB}{A.A.} & 12 & \textcolor{gruppeB}{B (B1)} & 10 & Fleißig, unsicher \\
\hline
\textcolor{gruppeB}{A.S.} & 12 & \textcolor{gruppeB}{B (B1)} & 9 & Einziger Junge, unsicher \\
\hline
\end{tabular}

\vspace{1em}

\textbf{Hinweis:} In dieser Aufstellung sind nur 8 SuS aufgeführt; laut Kursbeschreibung sind es 10 SuS (5 A / 5 B). Die Zuordnung kann je nach aktueller Zusammensetzung angepasst werden.

\subsection{Besondere Situation: Einstimmungsstunde}

\begin{infobox}[Didaktische Entscheidung]
\textbf{Warum KEINE Generalprobe?}
\begin{itemize}
    \item Generalprobe = \glqq{}Stundenverdopplung\grqq{} -- langweilig und kontraproduktiv
    \item Nervosität würde steigen statt sinken
    \item Gefahr: Fehler in Generalprobe demotivieren
\end{itemize}

\textbf{Stattdessen: Aktivierendes Wiederholen}
\begin{itemize}
    \item Bekannte Inhalte in \textbf{neuer Kombination}
    \item Gamification-Elemente für Motivation
    \item Positives Erfolgserlebnis vor der Beobachtung
    \item Selbstvertrauen stärken: \glqq{}Das kann ich schon!\grqq{}
\end{itemize}
\end{infobox}

% ============================================================================
% 2. SACHANALYSE
% ============================================================================
\section{Sachanalyse}

\subsection{Wiederholungsinhalte (keine neuen Strukturen!)}

\textbf{Diese Stunde wiederholt ausschließlich bereits eingeführte Strukturen:}

\begin{itemize}
    \item \textbf{Grammatik (bekannt):}
    \begin{itemize}
        \item \textit{ser + Adjektiv} (Eigenschaften): \textit{Soy simpático.}
        \item \textit{estar + Adjektiv} (Zustände): \textit{Estoy cansada.}
        \item \textit{tener + años/hermanos/...}: \textit{Tengo 16 años.}
        \item \textit{hay que + Infinitiv}: \textit{Hay que ser paciente.}
        \item \textit{porque + Begründung}: \textit{...porque me gusta.}
    \end{itemize}

    \item \textbf{Wortschatz (bekannt):}
    \begin{itemize}
        \item Adjektive: \textit{simpático/-a, inteligente, trabajador/-a, paciente, creativo/-a, divertido/-a, responsable, organizado/-a}
        \item Familie: \textit{hermano/-a, padre, madre, abuelo/-a}
        \item Hobbys/Aktivitäten: \textit{jugar, leer, escuchar música, hacer deporte}
        \item Berufe: \textit{médico, profesor/-a, bombero/-a, veterinario/-a}
    \end{itemize}

    \item \textbf{Kommunikative Strukturen (bekannt):}
    \begin{itemize}
        \item Sich vorstellen: \textit{Me llamo... Soy de... Tengo... años.}
        \item Nach Eigenschaften fragen: \textit{¿Cómo eres? ¿Cómo es tu...?}
        \item Personen beschreiben: \textit{Es... y también...}
    \end{itemize}
\end{itemize}

\subsection{Neue Kombination der bekannten Elemente}

\textbf{Das Neue ist nicht der Inhalt, sondern die Anwendung:}

\begin{tabular}{|l|p{9cm}|}
\hline
\textbf{Element} & \textbf{Neue Kombination/Anwendung} \\
\hline
Quiz-Duell & Tempo-Wettbewerb aktiviert abrufbares Wissen \\
\hline
Mystery-Karten & Zufällige Prompts erfordern spontane Reaktion \\
\hline
Tandem-Challenge & Bekannte Strukturen in Wettbewerbsformat \\
\hline
Peer-Beschreibung & Mitschüler beschreiben statt Familienmitglieder \\
\hline
\end{tabular}

% ============================================================================
% 3. DIDAKTISCHE ANALYSE
% ============================================================================
\section{Didaktische Analyse}

\subsection{Stellung in der Sequenz}

Dies ist die \textbf{9. Doppelstunde} (Tag vor der Beobachtung).

\vspace{0.5em}
\textbf{Sequenzübersicht:}
\begin{center}
\begin{tabular}{|c|l|l|}
\hline
\textbf{DS} & \textbf{Thema} & \textbf{Fokus} \\
\hline
1--5 & Getrennte Progression & A: LJ1 / B: LJ2 \\
\hline
6 & Brückenthema & Erster Overlap \\
\hline
7--8 & Integration & Gemeinsame Aktivitäten \\
\hline
\rowcolor{lightgray} \textbf{9} & \textbf{Repaso y Preparación} & \textbf{Einstimmung (Mo)} \\
\hline
10 & \textbf{Beobachtungsstunde} & Di 03.02.2026 \\
\hline
\end{tabular}
\end{center}

\subsection{Legitimation}

\textbf{Pädagogische Begründung:}
\begin{itemize}
    \item \textbf{Prüfungspsychologie:} Aktivierung am Vortag reduziert Nervosität
    \item \textbf{Spacing Effect:} Verteiltes Üben verbessert Langzeitretention
    \item \textbf{Selbstwirksamkeit:} Erfolgserlebnisse stärken Selbstvertrauen
    \item \textbf{Routine:} Bekannte Formate vermitteln Sicherheit
\end{itemize}

\textbf{Warum Gamification?}
\begin{itemize}
    \item Spielerischer Wettbewerb lenkt von Nervosität ab
    \item Punkte-System motiviert ohne Leistungsdruck
    \item Teamformate stärken Gruppenzusammenhalt vor Beobachtung
\end{itemize}

% ============================================================================
% 4. STUNDENZIELE
% ============================================================================
\section{Stundenziele}

\subsection{Grobziel}
Die SuS \textbf{aktivieren und festigen} ihre sprachlichen Mittel zur Personenbeschreibung, indem sie diese in spielerischen Formaten anwenden, sodass sie am Folgetag (Beobachtungsstunde) sicher und routiniert agieren können.

\subsection{Feinziele (operationalisiert)}

\begin{tabular}{|c|p{6.5cm}|c|c|c|}
\hline
\textbf{Nr.} & \textbf{Die SuS können...} & \textbf{AFB} & \textbf{Gruppe} & \textbf{Phase} \\
\hline
FZ1 & Vokabeln und Strukturen unter Zeitdruck abrufen & I & alle & Quiz \\
\hline
FZ2 & spontan auf unbekannte Prompts reagieren & II & alle & Speed-Dating \\
\hline
FZ3 & einen Mitschüler mit mind. 4 Sätzen beschreiben & II & alle & Tandem \\
\hline
FZ4 & Feedback zu Partnerbeschreibungen geben & II & alle & Tandem \\
\hline
FZ5 & ihre Stärken und Übungsfelder benennen & III & alle & Cool-Down \\
\hline
\end{tabular}

\vspace{1em}
\textbf{Hinweis:} Alle Ziele betreffen \textbf{bekannte Inhalte} -- es geht um Aktivierung, nicht um Neulernen!

% ============================================================================
% 5. VERLAUFSPLANUNG (PRÄSENZ)
% ============================================================================
\section{Verlaufsplanung (Präsenz-Variante)}

\textbf{Übersicht Doppelstunde (90 Min.)}

\begin{longtable}{|p{0.8cm}|p{1.5cm}|p{4cm}|p{3cm}|p{1.5cm}|p{2cm}|}
\hline
\textbf{Zeit} & \textbf{Phase} & \textbf{Lehrerhandeln} & \textbf{Schülerhandeln} & \textbf{SF} & \textbf{Medien} \\
\hline
\endhead

\multicolumn{6}{|c|}{\textbf{PHASE 1: REPASO RELÁMPAGO -- Quiz-Duell (15 Min.)}} \\
\hline

0--2 & Einstieg & Begrüßung: \glqq{}¡Hola! Hoy: Repaso divertido.\grqq{} Ankündigung Quiz & Zuhören, sich sammeln & Plenum & -- \\
\hline

2--5 & Vorbereitung & Teams einteilen (2 Teams à 5), Quiz-Regeln erklären & Teams bilden, aufstellen & Plenum & Tafel \\
\hline

5--14 & Quiz-Duell & Fragen vorlesen, Buzzer-System (Hand heben), Punkte notieren & Schnell antworten, jubeln & Teams & Quiz-Karten, Tafel \\
\hline

14--15 & Auswertung & Siegerteam verkünden, kurzes Lob & Klatschen, freuen & Plenum & -- \\
\hline

\multicolumn{6}{|c|}{\textbf{PHASE 2: LERNPFAD -- Differenzierte Übung (25 Min.)}} \\
\hline

15--17 & Überleitung & Lernpfad-Zugang erklären, Geräte verteilen & Tablets/Laptops holen & Plenum & Geräte \\
\hline

17--40 & Lernpfad & Beobachten, bei Bedarf helfen, individuelle Tipps & Selbstständig am Lernpfad arbeiten (A2 oder B1 Niveau) & EA & Lernpfad (digital) \\
\hline

40--42 & Abschluss LP & Kurze Rückmeldung: \glqq{}Wie weit seid ihr?\grqq{} & Fortschritt melden & Plenum & -- \\
\hline

\multicolumn{6}{|c|}{\textbf{PHASE 3: SPEED-DATING REMIX -- Mystery-Karten (20 Min.)}} \\
\hline

42--45 & Vorbereitung & Speed-Dating erklären, Stühle aufstellen (2 Reihen), Mystery-Karten verteilen & Plätze einnehmen, Karten sichten & Plenum & Stühle, Karten \\
\hline

45--60 & Speed-Dating & Timer stellen (2 Min. pro Runde), Gong für Wechsel, beobachten & Partner befragen nach Mystery-Prompt, dann wechseln & PA (rotierend) & Timer, Mystery-Karten \\
\hline

60--62 & Auswertung & \glqq{}Was war die überraschendste Antwort?\grqq{} & Kurze Berichte & Plenum & -- \\
\hline

\multicolumn{6}{|c|}{\textbf{PHASE 4: TANDEM-CHALLENGE -- Partnerbeschreibung (20 Min.)}} \\
\hline

62--65 & Vorbereitung & Teams bilden (A+B gemischt), Challenge erklären: \glqq{}Beschreibe deinen Partner!\grqq{} & Tandems finden, Material holen & Plenum & Notizzettel \\
\hline

65--72 & Interview & Partner interviewen (3 Fragen), Notizen machen & Gegenseitig befragen, notieren & Tandem & Notizzettel \\
\hline

72--80 & Beschreibung & Jedes Tandem: Eine Person beschreibt Partner, andere korrigiert/ergänzt & Beschreiben, Feedback geben & Tandem & -- \\
\hline

80--82 & Punkte & Beste Beschreibungen (flüssig, korrekt, kreativ) bekommen Bonuspunkte & Zuhören, bewerten (Daumen) & Plenum & Punkteliste \\
\hline

\multicolumn{6}{|c|}{\textbf{PHASE 5: COOL-DOWN + AUSBLICK (10 Min.)}} \\
\hline

82--85 & Reflexion & \glqq{}Was kannst du schon gut? Was übst du noch?\grqq{} & Kurze Selbsteinschätzung (Checkliste) & EA & Checkliste \\
\hline

85--88 & Ausblick & \glqq{}Morgen kommt jemand zu Besuch -- aber wir machen einfach weiter wie immer. Ihr könnt das!\grqq{} & Zuhören, nicken & Plenum & -- \\
\hline

88--90 & Abschluss & Positives Feedback, \glqq{}¡Hasta mañana!\grqq{} & Aufräumen, verabschieden & Plenum & -- \\
\hline

\end{longtable}

\textbf{Legende:} EA = Einzelarbeit, PA = Partnerarbeit, SF = Sozialform

% ============================================================================
% 6. DETAILPLANUNG DER GAMIFICATION-ELEMENTE
% ============================================================================
\section{Detailplanung der Gamification-Elemente}

\subsection{Phase 1: Repaso Relámpago (Quiz-Duell)}

\begin{gamebox}
\textbf{Format:} Team-Quiz mit Buzzer-System (Hand heben)

\textbf{Teams:}
\begin{itemize}
    \item Team 1: 2--3 aus Gruppe A + 2--3 aus Gruppe B (gemischt)
    \item Team 2: entsprechend
\end{itemize}

\textbf{Regeln:}
\begin{itemize}
    \item 10 Fragen (siehe Anlage)
    \item Wer zuerst die Hand hebt, darf antworten
    \item Richtig = 1 Punkt für Team
    \item Falsch = anderes Team darf
    \item Auf Spanisch antworten!
\end{itemize}

\textbf{Fragen-Typen:}
\begin{itemize}
    \item Vokabel-Blitz: \glqq{}Wie heißt \,Bruder' auf Spanisch?\grqq{}
    \item Grammatik-Check: \glqq{}\textit{Soy} oder \textit{estoy} cansado?\grqq{}
    \item Satz vervollständigen: \glqq{}Para ser médico hay que ser...\grqq{}
    \item Übersetzung: \glqq{}Ich bin 16 Jahre alt = ?\grqq{}
\end{itemize}

\textbf{Warum nicht langweilig?}
\begin{itemize}
    \item Wettbewerb erzeugt Spannung
    \item Tempo verhindert Grübeln
    \item Teamformat = sozialer Druck + Unterstützung
    \item Jubeln erlaubt!
\end{itemize}
\end{gamebox}

\subsection{Phase 3: Speed-Dating Remix (Mystery-Karten)}

\begin{gamebox}
\textbf{Format:} Rotierendes Partnerinterview mit Zufalls-Prompts

\textbf{Aufbau:}
\begin{itemize}
    \item 2 Stuhlreihen gegenüber (5 + 5)
    \item Eine Reihe rotiert nach 2 Minuten
    \item Ca. 5--6 Runden möglich
\end{itemize}

\textbf{Mystery-Karten:}

Jeder SuS zieht zu Beginn 3 Karten und muss in jeder Runde EINE verwenden. Die Karte gibt den Gesprächsimpuls vor.

\textit{Beispiele:}
\begin{itemize}
    \item \glqq{}Pregunta: ¿Cómo es tu mejor amigo/-a?\grqq{}
    \item \glqq{}Describe tu familia en 3 frases.\grqq{}
    \item \glqq{}¿Qué hay que ser para ser un buen estudiante?\grqq{}
    \item \glqq{}Di 5 adjetivos para describir a tu profesor/-a ideal.\grqq{}
    \item \glqq{}¿Cómo estás hoy? ¿Por qué?\grqq{}
\end{itemize}

\textbf{Warum nicht langweilig?}
\begin{itemize}
    \item Überraschungseffekt (welche Karte?)
    \item Neue Partner jede Runde
    \item Zeitdruck (nur 2 Min.)
    \item Bekanntes Format, aber unvorhersehbare Prompts
\end{itemize}
\end{gamebox}

\subsection{Phase 4: Tandem-Challenge}

\begin{gamebox}
\textbf{Format:} A+B gemischt, gegenseitige Beschreibung mit Punkte-Wertung

\textbf{Tandem-Zuordnung (Vorschlag):}
\begin{tabular}{ll}
\textcolor{gruppeA}{E.H.} + \textcolor{gruppeB}{A.S.} & (TOP + unsicher -- E.H. als Vorbild) \\
\textcolor{gruppeA}{L.N.} + \textcolor{gruppeB}{H.P.} & (proaktiv + Nachzügler) \\
\textcolor{gruppeA}{C.P.} + \textcolor{gruppeB}{A.A.} & (ruhig + unsicher -- beides ruhig) \\
\textcolor{gruppeA}{V.S.} + \textcolor{gruppeA}{L.K.} & (fleißig + abgelenkt) \\
Ggf. 5. Tandem & je nach Anwesenheit \\
\end{tabular}

\textbf{Ablauf:}
\begin{enumerate}
    \item Interview (3 Min.): 3 Fragen an Partner
    \item Beschreibung schreiben (4 Min.): 4 Sätze über Partner
    \item Vorlesen (3 Min.): Eine Person liest vor
    \item Feedback (2 Min.): Partner korrigiert/ergänzt
\end{enumerate}

\textbf{Punkte-System:}
\begin{itemize}
    \item 1 Punkt: Beschreibung verständlich
    \item 2 Punkte: Mind. 4 Sätze, grammatisch weitgehend korrekt
    \item 3 Punkte: Flüssig, kreativ, mit Konnektoren
    \item Bonus: Klasse gibt Daumen-Feedback
\end{itemize}

\textbf{Warum nicht langweilig?}
\begin{itemize}
    \item Peer-Fokus (über Mitschüler reden = interessant)
    \item Wettbewerbsformat
    \item Sofortiges Feedback vom Beschriebenen
\end{itemize}
\end{gamebox}

% ============================================================================
% 7. ONLINE-VARIANTE
% ============================================================================
\newpage
\section{Online-Variante (bei Lehrerausfall)}

\begin{onlinebox}
Falls die Lehrkraft am 02.02.2026 ausfällt, bearbeiten die SuS selbstständig folgende Aufgaben über itslearning:

\vspace{1em}
\textbf{Zeitrahmen: 90 Minuten (selbstorganisiert)}

\begin{tabular}{|c|p{5cm}|c|p{5cm}|}
\hline
\textbf{Teil} & \textbf{Aktivität} & \textbf{Zeit} & \textbf{Abgabe} \\
\hline
1 & Lernpfad (LP-09) & 45 Min. & Automatisch gespeichert \\
\hline
2 & Schreibaufgabe (siehe unten) & 30 Min. & itslearning (Textfeld) \\
\hline
3 & Selbst-Checkliste ausfüllen & 15 Min. & itslearning (Formular) \\
\hline
\end{tabular}

\vspace{1em}
\textbf{Teil 1: Lernpfad (45 Min.)}
\begin{itemize}
    \item URL: [Link zu LP-09]
    \item Differenziert: A2-Track oder B1-Track
    \item Automatische Korrektur und Feedback
    \item Fortschritt wird gespeichert (localStorage)
\end{itemize}

\vspace{1em}
\textbf{Teil 2: Schreibaufgabe (30 Min.)}

\textit{Aufgabe:} \glqq{}Describe a tu compañero/-a ideal para trabajar en clase.\grqq{}

\begin{itemize}
    \item Mindestens 5 Sätze
    \item Verwende: \textit{ser + Adjektiv}, \textit{también}, \textit{porque}
    \item Beschreibe Eigenschaften, die für Gruppenarbeit wichtig sind
\end{itemize}

\textit{Hilfe:}
\begin{itemize}
    \item \textit{Mi compañero/-a ideal es...}
    \item \textit{También es... porque...}
    \item \textit{No es... sino...}
\end{itemize}

\vspace{1em}
\textbf{Teil 3: Selbst-Checkliste (15 Min.)}

Die SuS füllen folgendes Formular aus:

\begin{tabular}{|p{6cm}|c|c|c|}
\hline
\textbf{Das kann ich schon:} & \textbf{Sicher} & \textbf{Geht so} & \textbf{Übe ich noch} \\
\hline
Mich vorstellen (Name, Alter, Herkunft) & $\Box$ & $\Box$ & $\Box$ \\
\hline
Adjektive mit \textit{ser} verwenden & $\Box$ & $\Box$ & $\Box$ \\
\hline
\textit{ser} und \textit{estar} unterscheiden & $\Box$ & $\Box$ & $\Box$ \\
\hline
Personen mit 3+ Adjektiven beschreiben & $\Box$ & $\Box$ & $\Box$ \\
\hline
Konnektoren verwenden (\textit{también, pero, porque}) & $\Box$ & $\Box$ & $\Box$ \\
\hline
Fragen stellen (\textit{¿Cómo eres?}) & $\Box$ & $\Box$ & $\Box$ \\
\hline
\textit{hay que ser + Adjektiv} verwenden & $\Box$ & $\Box$ & $\Box$ \\
\hline
\end{tabular}

\vspace{1em}
\textbf{Hinweis an SuS:}

\glqq{}Morgen ist die Beobachtungsstunde. Diese Übungen helfen dir, dich vorzubereiten. Du hast alles schon gelernt -- jetzt geht es nur darum, es zu aktivieren. ¡Tú puedes!\grqq{}

\end{onlinebox}

% ============================================================================
% 8. ERWARTETE SCHWIERIGKEITEN
% ============================================================================
\section{Erwartete Schwierigkeiten}

\begin{tabular}{|p{4.5cm}|p{8.5cm}|}
\hline
\textbf{Schwierigkeit} & \textbf{Reaktion} \\
\hline
\textbf{SuS nervös wegen Beobachtung} & Explizit thematisieren: \glqq{}Morgen ist ein ganz normaler Tag.\grqq{} Fokus auf Spaß, nicht auf Perfektion. \\
\hline
\textbf{Quiz zu schnell für schwächere SuS} & Einfache Fragen gezielt an schwächere SuS richten (vorher markieren). Team-Hilfe erlauben. \\
\hline
\textbf{Speed-Dating: SuS wissen nichts zu sagen} & Mystery-Karten geben klare Prompts. L geht herum und hilft. \\
\hline
\textbf{Tandem-Challenge: Ungleiche Paare} & Paare bewusst so zusammenstellen, dass stärkere SuS unterstützen können. \\
\hline
\textbf{Zeitproblem (Aktivitäten zu lang)} & Puffer in Phase 5 (Cool-Down kann gekürzt werden). Speed-Dating auf 4 statt 6 Runden. \\
\hline
\textbf{Motivationsproblem (\glqq{}Ist doch nur Wiederholung\grqq{})} & Gamification-Punkte sammeln. Ankündigung: \glqq{}Bestes Team bekommt [kleine Belohnung].\grqq{} \\
\hline
\textbf{Technische Probleme beim Lernpfad} & Alternativ: AB-09 (analog) bereithalten mit gleichen Übungen. \\
\hline
\end{tabular}

% ============================================================================
% 9. HINWEISE ZUM AUSBLICK (PHASE 5)
% ============================================================================
\section{Hinweise zum Ausblick (Phase 5)}

\begin{infobox}[Formulierungshilfen für den Ausblick]
\textbf{RICHTIG (positiv, entlastend):}
\begin{itemize}
    \item \glqq{}Morgen kommt jemand zu Besuch, der sich für unseren Unterricht interessiert.\grqq{}
    \item \glqq{}Ihr habt heute gezeigt, dass ihr das alles könnt!\grqq{}
    \item \glqq{}Wir machen morgen einfach weiter wie immer.\grqq{}
    \item \glqq{}Es ist kein Test -- einfach normaler Unterricht.\grqq{}
    \item \glqq{}¡Muy bien! ¡Estoy seguro/-a de que mañana va a ser genial!\grqq{}
\end{itemize}

\textbf{FALSCH (Nervosität erzeugend):}
\begin{itemize}
    \item[\texttimes] \glqq{}Morgen ist die wichtige Beobachtungsstunde.\grqq{}
    \item[\texttimes] \glqq{}Bitte benehmt euch morgen gut.\grqq{}
    \item[\texttimes] \glqq{}Ich werde bewertet, also...\grqq{}
    \item[\texttimes] \glqq{}Zeigt morgen, was ihr könnt!\grqq{} (impliziert Leistungsdruck)
\end{itemize}
\end{infobox}

% ============================================================================
% 10. ANLAGEN
% ============================================================================
\section{Anlagen}

\begin{enumerate}
    \item Quiz-Fragen für Repaso Relámpago (10 Fragen)
    \item Mystery-Karten für Speed-Dating (20 Karten)
    \item Lernpfad-Struktur (LP-09)
    \item Selbst-Checkliste (Print-Version)
    \item Punkteliste für Gamification
\end{enumerate}

\subsection{Anlage 1: Quiz-Fragen für Repaso Relámpago}

\begin{tabular}{|c|p{7cm}|p{4.5cm}|}
\hline
\textbf{\#} & \textbf{Frage} & \textbf{Antwort} \\
\hline
1 & Wie heißt \glqq{}Bruder\grqq{} auf Spanisch? & \textit{hermano} \\
\hline
2 & \textit{Soy} oder \textit{estoy} cansado/-a? & \textit{estoy} (Zustand) \\
\hline
3 & Übersetze: \glqq{}Ich bin 16 Jahre alt.\grqq{} & \textit{Tengo 16 años.} \\
\hline
4 & Welches Adjektiv passt: nett, sympathisch = ? & \textit{simpático/-a} \\
\hline
5 & Vervollständige: Para ser médico hay que ser... & \textit{paciente / responsable / ...} \\
\hline
6 & Was bedeutet \textit{también}? & auch \\
\hline
7 & \textit{Soy} oder \textit{estoy} inteligente? & \textit{soy} (Eigenschaft) \\
\hline
8 & Wie fragt man: \glqq{}Wie bist du?\grqq{} & \textit{¿Cómo eres?} \\
\hline
9 & Nenne 3 Adjektive für eine gute Lehrerin. & z.B. \textit{paciente, inteligente, divertida} \\
\hline
10 & Übersetze: \glqq{}Mein Vater ist lustig.\grqq{} & \textit{Mi padre es divertido.} \\
\hline
\end{tabular}

\subsection{Anlage 2: Mystery-Karten (Auswahl)}

\begin{verbatim}
+--------------------------------------------+
|  MYSTERY CARD #1                           |
|  ----------------------------------------  |
|  Pregunta: ¿Cómo es tu mejor amigo/-a?     |
+--------------------------------------------+

+--------------------------------------------+
|  MYSTERY CARD #2                           |
|  ----------------------------------------  |
|  Describe tu familia en 3 frases.          |
+--------------------------------------------+

+--------------------------------------------+
|  MYSTERY CARD #3                           |
|  ----------------------------------------  |
|  ¿Qué hay que ser para ser un buen         |
|  estudiante? Nombra 3 características.     |
+--------------------------------------------+

+--------------------------------------------+
|  MYSTERY CARD #4                           |
|  ----------------------------------------  |
|  ¿Cómo estás hoy? ¿Por qué?                |
+--------------------------------------------+

+--------------------------------------------+
|  MYSTERY CARD #5                           |
|  ----------------------------------------  |
|  Di 5 adjetivos positivos sobre ti.        |
+--------------------------------------------+

+--------------------------------------------+
|  MYSTERY CARD #6                           |
|  ----------------------------------------  |
|  ¿Cómo es tu profesor/-a ideal?            |
+--------------------------------------------+

+--------------------------------------------+
|  MYSTERY CARD #7                           |
|  ----------------------------------------  |
|  Pregunta: ¿Tienes hermanos? ¿Cómo son?    |
+--------------------------------------------+

+--------------------------------------------+
|  MYSTERY CARD #8                           |
|  ----------------------------------------  |
|  ¿Qué características son importantes      |
|  para trabajar en grupo?                   |
+--------------------------------------------+
\end{verbatim}

\subsection{Anlage 3: Lernpfad-Struktur (LP-09)}

\begin{tabular}{|c|l|c|p{5cm}|}
\hline
\textbf{\#} & \textbf{Modul} & \textbf{Typ} & \textbf{Inhalt} \\
\hline
1 & Vokabel-Check & Exercise & 10 Adjektive zuordnen (Drag\&Drop) \\
\hline
2 & \textit{ser/estar}-Übung & Exercise & 8 Lückensätze \\
\hline
3 & Personenbeschreibung & Exercise & Bild beschreiben (Multiple Choice) \\
\hline
4 & \textit{hay que ser...} & Exercise & Berufe + Eigenschaften zuordnen \\
\hline
5 & Konnektoren & Exercise & Sätze verbinden \\
\hline
6 & Mini-Dialog & Exercise & Dialog vervollständigen \\
\hline
7 & Selbsttest & Test & 5 Fragen zur Selbsteinschätzung \\
\hline
\end{tabular}

\textbf{Differenzierung:}
\begin{itemize}
    \item \textcolor{gruppeA}{\textbf{A2-Track:}} Einfachere Vokabeln, mehr Scaffolding, Hilfe-Buttons
    \item \textcolor{gruppeB}{\textbf{B1-Track:}} Komplexere Sätze, weniger Hilfe, freiere Produktion
\end{itemize}

\subsection{Anlage 4: Selbst-Checkliste (Print-Version)}

\begin{verbatim}
+---------------------------------------------------------------------+
|  SELBST-CHECKLISTE: ¿Qué puedo ya?                                  |
|  Name: _________________________  Datum: 02.02.2026                 |
+---------------------------------------------------------------------+
|                                                                     |
|  Kreuze an, wie sicher du dich fühlst:                              |
|                                                                     |
|                                    SICHER  GEHT SO  ÜBE ICH NOCH   |
|  ------------------------------------------------------------------ |
|  Mich vorstellen                    [ ]       [ ]        [ ]        |
|  Adjektive mit ser verwenden        [ ]       [ ]        [ ]        |
|  ser und estar unterscheiden        [ ]       [ ]        [ ]        |
|  Personen beschreiben (3+ Sätze)    [ ]       [ ]        [ ]        |
|  Konnektoren verwenden              [ ]       [ ]        [ ]        |
|  Fragen stellen                     [ ]       [ ]        [ ]        |
|  hay que ser + Adjektiv             [ ]       [ ]        [ ]        |
|                                                                     |
|  ------------------------------------------------------------------ |
|  Das nehme ich mir für morgen vor:                                  |
|                                                                     |
|  __________________________________________________________________ |
|                                                                     |
|  __________________________________________________________________ |
|                                                                     |
+---------------------------------------------------------------------+
\end{verbatim}

% ============================================================================
% 11. DIDAKTIK-SCORE
% ============================================================================
\newpage
\section{Didaktik-Score}

\begin{scorebox}
\textbf{Bewertung der didaktischen Qualität nach 10 Dimensionen}\\
\textbf{Ziel-Score: $\geq$ 9.0 (Premium-Qualität)}

\vspace{1em}
\begin{tabular}{|c|l|c|c|p{4.5cm}|}
\hline
\textbf{\#} & \textbf{Dimension} & \textbf{Gew.} & \textbf{Score} & \textbf{Notiz} \\
\hline
1 & Differenzierung & 15\% & 9/10 & A/B-Tracks im Lernpfad, gemischte Tandems \\
\hline
2 & Taxonomie (AFB) & 10\% & 9/10 & I (Quiz), II (Beschreiben), III (Reflexion) \\
\hline
3 & Lernziel-Alignment & 10\% & 10/10 & Reine Wiederholung bekannter RP-Inhalte \\
\hline
4 & Aktivierung & 10\% & 10/10 & Durchgehend interaktiv, keine Passivphasen \\
\hline
5 & Scaffolding & 10\% & 10/10 & Mystery-Karten, Checkliste, Lernpfad-Hilfen \\
\hline
6 & Feedback-Qualität & 10\% & 9/10 & Auto-Korrektur LP, Peer-Feedback Tandem \\
\hline
7 & Sprachliche Klarheit & 10\% & 10/10 & Bekannte Strukturen, keine neuen Begriffe \\
\hline
8 & Motivation \& Relevanz & 10\% & 10/10 & Gamification, Punkte, Wettbewerb, Spaß \\
\hline
9 & Inklusion (UDL) & 10\% & 9/10 & Digital + analog, mündlich + schriftlich \\
\hline
10 & Praktikabilität & 5\% & 10/10 & 90 Min. realistisch, Material einfach \\
\hline
\multicolumn{3}{|r|}{\textbf{GESAMT:}} & \textbf{9.5/10} & \\
\hline
\end{tabular}

\vspace{1em}
\textbf{Bewertungsschwellen:}
\begin{itemize}
    \item[\checkmark] \textcolor{successgreen}{$\geq$ 9.0: \textbf{FREIGABE} (Premium-Qualität)}
    \item[$\Box$] \textcolor{warningorange}{8.0--8.9: Iteration empfohlen}
    \item[$\Box$] 6.0--7.9: Überarbeitung erforderlich
    \item[$\Box$] $<$ 6.0: Grundlegende Neukonzeption
\end{itemize}

\vspace{1em}
\textcolor{successgreen}{\textbf{STATUS: FREIGABE}}

\vspace{1em}
\textbf{Stärken:}
\begin{itemize}
    \item Klares Konzept: Wiederholung OHNE Langeweile durch Gamification
    \item Keine Generalprobe -- psychologisch klug
    \item Online-Variante als Backup vollständig durchdacht
    \item Ausblick-Formulierungen explizit positiv
    \item Alle Aktivitäten nutzen bekannte Inhalte in neuer Form
\end{itemize}

\textbf{Optionale Verbesserungen:}
\begin{itemize}
    \item Kahoot statt analogem Quiz (falls technisch möglich)
    \item Kleine Preise für Siegerteam vorbereiten (Aufkleber, Süßigkeit)
\end{itemize}

\end{scorebox}

\end{document}
